% Generated draft: requires mathematical and source review.
\documentclass[11pt]{article}
\usepackage[margin=25mm]{geometry}
\usepackage{fontspec}
\setmainfont{Latin Modern Roman}
\usepackage{amsmath,amssymb,mathtools,unicode-math}
\setmathfont{Latin Modern Math}
\usepackage{enumitem,needspace,xurl,hyperref}
\setlength{\parindent}{0pt}
\setlength{\parskip}{5pt}
\setlength{\emergencystretch}{3em}
\newcommand{\N}{\mathbb{N}}
\newcommand{\Z}{\mathbb{Z}}
\newcommand{\Q}{\mathbb{Q}}
\newcommand{\R}{\mathbb{R}}
\newcommand{\C}{\mathbb{C}}
\newcommand{\F}{\mathbb{F}}
\newcommand{\A}{\mathbb{A}}
\newcommand{\PP}{\mathbb{P}}
\newcommand{\E}{\mathbb{E}}
\newcommand{\eps}{\varepsilon}
\newcommand{\dd}{\,\mathrm{d}}
\begin{document}
\section*{UnsolvedMath ID: 6700058}
\textbf{Source problem code:} catalogue entry
\subsection{Definitions and mathematical statement}
Interpret the source's infinite-dimensional Hilbert space \(\mathbb R^\infty\) as a real Hilbert space \(H\), conventionally \(\ell^2\), and let \(B_H(0,1)\) be its open unit ball. The source speculates that every manifold might admit an immersion into this ball with principal curvatures bounded by one universal constant, suggesting \(10^6\).

For a \(C^2\) immersion \(F:M\to H\), use its induced metric and second fundamental form \(\mathrm{II}_F\). In higher codimension a precise absolute principal-curvature bound can be expressed by
\[
 \|A_\nu\|_{\mathrm{op}}\le K
 \quad\text{for every unit normal }\nu,
 \qquad
 \langle A_\nu v,w\rangle=\langle\mathrm{II}_F(v,w),\nu\rangle.
\]
With this convention the proposed universal assertion is
\[
 \exists K<\infty\ \forall M\text{ in the intended smooth-manifold class}\ 
 \exists F:M\looparrowright B_H(0,1):
 \quad\sup_{x,\,|\nu|=1}\|A_\nu(x)\|_{\mathrm{op}}\le K,
\]
and the suggested numerical version takes \(K=1{,}000{,}000\). The constant is independent of the manifold and, on the literal ``all manifolds'' reading, its dimension.

The source does not specify manifold dimension range, compactness, boundary, completeness of the induced metric, or a prescribed metric to be preserved. Thus no isometric-immersion or embedding requirement is added. Its higher-codimension curvature convention and the meaning of \(\mathbb R^\infty\) also need confirmation. The research summary sometimes uses embedding and finite-codimension language, which is not substituted for the statement's immersion into a Hilbert ball.

The wording is a speculative possibility, not an asserted theorem or a claimed sharp value of \(K\). The summary reports related finite-dimensional and curvature-bounded work and mentions Nash's \(C^1\) embeddings, without exact references; \(C^1\) existence by itself does not specify a classical second-fundamental-form bound. The reported open status is source information without independent current verification.
\subsection{Short English statement}
Could every manifold be immersed in a Hilbert-space unit ball with one universal bound on all normal-direction principal curvatures, for example \(10^6\)? The source asks about immersions and does not prescribe an isometry, embedding, or completeness requirement.
\subsection{Sources}
\begin{itemize}
\item Local registry: \url{data/unsolvedmath/problems.json}, permanent UnsolvedMath ID \(6700058\), source code catalogue entry; supplied statement, background, and research summary, updated 2026-08-06.
\item Misha Gromov, \emph{101 Questions, Problems and Conjectures around Scalar Curvature} (2017), printed tag [?62], occurrence 61, PDF page 59: \url{https://www.ihes.fr/~gromov/wp-content/uploads/2018/08/101-problemsOct1-2017.pdf}. Supplied normalized PDF-text extraction accessed 2026-07-29; described as incomplete.
\item The summary mentions Nash and related curvature-bounded immersion/embedding work without separate publication details or a precise theorem for this Hilbert-ball assertion.
\end{itemize}
\end{document}
